\documentclass[12pt,a5paper]{article}

% Russian language
\usepackage[utf8]{inputenc}
\usepackage[T2A]{fontenc}
\usepackage[russian]{babel}

% Page layout
\usepackage[a5paper,margin=1.5cm]{geometry}
\usepackage{setspace}
\setstretch{1.08}

% Graphics and figures
\usepackage{graphicx}
\usepackage{float}

% Tables
\usepackage{longtable}
\usepackage{booktabs}
\usepackage{array}
\usepackage{tabularx}

% Mathematics
\usepackage{amsmath}
\usepackage{amssymb}
\usepackage{amsthm}

% Hyperlinks
\usepackage{hyperref}
\hypersetup{
    colorlinks=true,
    linkcolor=black,
    citecolor=black,
    urlcolor=black
}

% Headers and footers
\usepackage{fancyhdr}
\pagestyle{fancy}
\fancyhf{}
\renewcommand{\headrulewidth}{0pt}
\fancyfoot[C]{\thepage}

% Section formatting
\usepackage{titlesec}
\titleformat{\section}{\centering\normalsize\bfseries}{\thesection}{1em}{}
\titleformat{\subsection}{\normalsize\bfseries}{\thesubsection}{1em}{}

% Reduce spacing
\titlespacing*{\section}{0pt}{8pt}{4pt}
\titlespacing*{\subsection}{0pt}{6pt}{3pt}

% Paragraph settings
\setlength{\parindent}{1.25cm}
\setlength{\parskip}{0pt}

\begin{document}
\sloppy

% Suppress page number on first page
\thispagestyle{empty}

% ============================================================
% TITLE PAGE
% ============================================================
\begin{center}
\textbf{Национальный исследовательский ядерный университет <<МИФИ>>}

\vspace{1cm}

\textit{На правах рукописи}

\vspace{2cm}

\textbf{Анаэдевха Роджер Ник}

\vspace{1.5cm}

{\large \textbf{Методы и алгоритмы повышения
робастности, устойчивости, эффективности и точности динамических
графовых нейронных сетей в состязательных,
распределённых и нестационарных условиях}}

\vspace{1.5cm}

\textbf{2.3.1.\ Системный анализ, управление и обработка\\
информации, статистика\\
(технические науки)}

\vspace{2cm}

\textbf{АВТОРЕФЕРАТ}

\vspace{0.5cm}

диссертации на соискание учёной степени\\
кандидата технических наук

\vfill

Москва -- 2026
\end{center}

\newpage

% ============================================================
% PAGE 2 — Committee and defence details
% ============================================================
\thispagestyle{empty}
\enlargethispage{1.5cm}

{\scriptsize
\begin{center}
\textbf{Работа выполнена в федеральном государственном автономном
образовательном учреждении высшего образования
<<Национальный исследовательский ядерный университет <<МИФИ>>}
\end{center}

\vspace{0.15cm}

\noindent\begin{tabular}{@{}p{2.8cm}p{7.4cm}@{}}
\textbf{Научный руководитель:} & \textbf{Трофимов Александр Геннадьевич} \newline
доктор физико-математических наук, профессор, профессор кафедры №\,22 <<Кибернетика>> Национального исследовательского ядерного университета <<МИФИ>> \\
\end{tabular}

\vspace{0.15cm}

\noindent\textbf{Официальные оппоненты:}
\vspace{0.03cm}

\noindent\begin{tabular}{@{}p{2.8cm}p{7.4cm}@{}}
доктор технических наук, профессор, лауреат Премии Президента РФ в области образования, почётный работник науки и высоких технологий РФ
\newline \textbf{Еремеев Александр Павлович} &
профессор кафедры прикладной математики и искусственного интеллекта ФГБОУ ВО <<Национальный исследовательский университет <<МЭИ>> \\[0.1cm]

доктор технических наук, профессор
\newline \textbf{Калянов Георгий Николаевич} &
главный научный сотрудник лаборатории инфраструктурных систем ФГБУ Институт проблем управления им.~В.\,А.\,Трапезникова Российской академии наук \\
\end{tabular}

\vspace{0.15cm}

\noindent\begin{tabular}{@{}p{2.8cm}p{7.4cm}@{}}
\textbf{Ведущая организация:} & Федеральное государственное бюджетное образовательное учреждение высшего образования <<Тверской государственный технический университет>> \\
\end{tabular}

\vspace{0.15cm}

\noindent Защита состоится 30.06.2026 на заседании диссертационного совета~2.01 федерального государственного автономного образовательного учреждения высшего образования <<Национальный исследовательский ядерный университет <<МИФИ>> (115409, г.~Москва, Каширское шоссе, 31)

\vspace{0.1cm}

\noindent С диссертацией можно ознакомиться в библиотеке и на сайте \url{https://ds.mephi.ru/} федерального государственного автономного образовательного учреждения высшего образования <<Национальный исследовательский ядерный университет <<МИФИ>>

\vspace{0.1cm}

\noindent Автореферат разослан \underline{\hspace{2.5cm}} 2026 года

\vspace{0.2cm}

\noindent\begin{tabular}{@{}p{4.2cm}p{2cm}p{3.5cm}@{}}
Учёный секретарь диссертационного совета, кандидат физико-математических наук & & Климанов Сергей Геннадьевич \\
\end{tabular}
}

\newpage
% ============================================================
% MAIN CONTENT BEGINS
% ============================================================
\setcounter{page}{3}

\section*{ОБЩАЯ ХАРАКТЕРИСТИКА РАБОТЫ}
\addcontentsline{toc}{section}{Общая характеристика работы}

\subsection*{Актуальность темы исследования и степень её разработанности}

Графовые нейронные сети стали основной парадигмой глубокого обучения для извлечения предсказательных представлений из реляционных данных, однако их механизм передачи сообщений делает их чувствительными к возмущениям топологии графа, признаков узлов и временного порядка событий. Эта чувствительность имеет точное происхождение в двойственной природе входов графовых нейронных сетей: в отличие от обычных нейронных сетей, принимающих единый тензор данных, графовая нейронная сеть одновременно получает два входа --- числовые \emph{входные данные} ($\mathbf{X}\in\mathbb{R}^{n\times d}$, матрица признаков, кодирующая атрибуты сущностей) и реляционную \emph{входную структуру} ($\mathbf{A}$, матрица смежности, кодирующая попарные связи) --- и формирует выход путём распространения данных \emph{вдоль} структуры через итеративную агрегацию окрестностей (рис.~\ref{fig:dual_input_ar}). Возмущение любого из входов может нарушить предсказания. В настоящей диссертации исследуется робастность обученных представлений на динамических графовых структурах как фундаментальная задача прикладной математики и искусственного интеллекта.

\begin{figure}[H]
\centering
\includegraphics[width=0.89\textwidth]{figures/fig1a_gnn.pdf}
\vspace{-0.5cm}
\caption{Двойственная природа входов графовых нейронных сетей: данные ($\mathbf{X}$) и структура ($\mathbf{A}$) обрабатываются совместно; возмущение любого из них нарушает обученное представление}
\label{fig:dual_input_ar}
\end{figure}

Большинство существующих графовых нейронных сетей работают в дискретном времени, принудительно привязывая нерегулярные события графа к фиксированной тактовой сетке и теряя как информацию, так и гарантии робастности. Формулировки непрерывного времени на основе нейронных обыкновенных дифференциальных уравнений гладко следуют за временем событий, захватывают множественные временные масштабы и обнаруживают границу Липшица, дающую сертифицированную робастность --- преимущество, которое ни одна дискретная модель не может обеспечить. Неравенство Грёнуолла формализует это: $\|\mathbf{h}_1(T)-\mathbf{h}_2(T)\|\le e^{L_g T}\|\mathbf{x}_1-\mathbf{x}_2\|$. Однако модели непрерывного времени получают ограниченное внимание в литературе. Распределённая среда --- в которой несколько организаций владеют частью одних и тех же графовых данных и не могут ими делиться --- требует федеративного обучения с формальными гарантиями приватности. Существующие модели на основе трансформеров масштабируются квадратично ($\mathcal{O}(L^2)$), но структурированные модели пространства состояний снижают это до $\mathcal{O}(L\log L)$. Кроме того, реальные графовые данные дрейфуют после развёртывания, и статические модели незаметно теряют точность без предупреждения.

Четыре взаимосвязанные открытые проблемы концентрируют нерешённые задачи: (1)~динамика непрерывного времени на эволюционирующих графах с сертифицированными гарантиями робастности; (2)~федеративное обучение на распределённых фрагментах графов с одновременной византийской устойчивостью, дифференциальной приватностью и выравниванием распределений; (3)~вычислительная эффективность для крупномасштабных темпоральных графов; и (4)~онлайн-адаптация в условиях нестационарности с калиброванной оценкой неопределённости. Эти проблемы образуют логическую цепочку: динамика непрерывного времени обеспечивает основу представлений, федеративные алгоритмы обеспечивают распределённое развёртывание, эффективный вывод позволяет работу в реальном времени, а онлайн-адаптация гарантирует долгосрочную надёжность.

Значительный вклад внесли Kipf, Welling, Veli\v{c}kovi\'{c}, Hamilton (графовые нейронные сети), Chen (нейронные ОДУ), Madry, Carlini, Wagner (состязательное МО), McMahan (федеративное обучение) и Gu, Dao (модели пространства состояний). Однако ни один существующий подход не решает совместно все четыре требования при трёх рабочих условиях.

Графовые нейронные сети внедряются в критически важные системы --- здравоохранение, финансы, сетевую безопасность, автономные системы --- быстрее, чем могут быть установлены их гарантии. Во всех случаях защитные механизмы должны быть встроены в архитектуру модели. Эти четыре проблемы отображаются на три рабочих условия и четыре сквозных требования, пересечения которых определяют архитектурную основу данной диссертации.

Три условия не являются взаимоисключающими: реальные системы сталкиваются с двумя или всеми тремя одновременно. На рис.~\ref{fig:venn_ar} показано, как семь методов, разработанных в данной диссертации (М1--М7), расположены в пространстве пересекающихся условий, при этом М4: MambaShield находится в центре всех трёх, а М1: CT-TGNN связывает состязательные и нестационарные условия через сертифицированную динамику непрерывного времени.

\begin{figure}[!t]
\centering
\includegraphics[width=0.89\columnwidth]{figures/fig2a_t_by_f.pdf}
\vspace{-0.5cm}
\caption{Три условия (слева), четыре требования (справа) и их взаимосвязь. \textbf{Примечание:} Только модели типа MambaShield ближе всего к выполнению всех 3 условий ГНС. Также ни одно требование не может иметь все 3 процесса как <<Полная обработка графа>>, а должно иметь по крайней мере 1 как <<Частичная обработка графа>>, чтобы остальные 2 были <<Полная обработка графа>>.}
\label{fig:venn_ar}
\end{figure}

На рис.~\ref{fig:main_arch} данное отображение формализуется как архитектурная основа исследовательской программы, показывающая три рабочих условия (состязательная среда, нестационарность, распределённая динамика) $\times$ четыре требования (робастность, точность, эффективность, приватность). Центральная матрица сопоставляет каждый из семи методов пересечениям условие--требование, покрывая все двенадцать ячеек.

\begin{figure}[!t]
\centering
\includegraphics[width=0.96\textwidth]{figures/Main_arch_v1.pdf}
\vspace{0.5cm}
\caption{Архитектурная основа исследовательской программы}
\label{fig:main_arch}
\end{figure}

Определив пространство задач, условия функционирования решений и архитектурную основу, организующую исследование, в следующих разделах определяются объект и предмет исследования, формулируются цель и задачи работы, описывается научная новизна разработанных методов.

\subsection*{Объект и предмет исследования}

\textbf{Объектом исследования} являются динамические графовые нейронные сети (ГНС), функционирующие в состязательных, нестационарных и распределённых средах.

\textbf{Предметом исследования} являются методы обеспечения робастности, эффективности, дифференциальной приватности и точности динамических ГНС, функционирующих в состязательных, нестационарных и распределённых средах.

\subsection*{Цель и задачи исследования}

\textbf{Целью} настоящей диссертации является повышение робастности графовых нейронных сетей к состязательным возмущениям, их устойчивости в нестационарных условиях, эффективности и точности для использования в реальном времени, а также обеспечение приватности в распределённых средах.

\textbf{Задачи:}

\begin{enumerate}
\item Разработать модели непрерывного времени (НВ) на основе нейронных обыкновенных дифференциальных уравнений с многомасштабной темпоральной динамикой и сертифицированной робастностью по Липшицу.

\item Разработать методы, позволяющие НВ-моделям функционировать в распределённых средах на основе федеративного обучения с византийско-устойчивой агрегацией, дифференциальной приватностью и выравниванием распределений методом оптимального транспорта.

\item Разработать методы снижения стоимости компромиссов и повышения эффективности как монолитных, так и распределённых моделей на основе многоуровневых вложений в сочетании с селективными архитектурами пространства состояний (MambaShield), снижающими сложность с квадратичной до близкой к линейной.

\item Разработать методы оценки и адаптации ГНС к нестационарным средам --- то есть к дрейфу распределений --- на основе иерархических гауссовских процессов для калиброванной оценки неопределённости и онлайновых байесовских апостериорных обновлений.

\item Интегрировать все методы в единую систему, провести валидацию на крупномасштабных бенчмарках и доказать её качество для готовности к развёртыванию с помощью регуляризации согласованности и комплексного тестирования.

\item Адаптировать единую систему для задач кибербезопасности, учитывая специфику данной области применения, с помощью предметно-ориентированных адаптеров и уточнения модели угроз.

\item Применить единую систему и продемонстрировать её качество в области гибридных систем обнаружения и предотвращения сетевых вторжений (СОПСВ) путём развёртывания в реальном времени на программной платформе RobustIDPS.ai.
\end{enumerate}

\subsection*{Методология и методы исследования}

В работе использованы методы системного анализа, теории графов, дифференциальных уравнений, функционального анализа, теории оптимального транспорта, глубокого обучения, байесовского вывода, теории информации, теории игр, спектральной теории графов, теории оптимизации, теории вероятностей, математической статистики и вычислительного эксперимента.

\subsection*{Научная новизна работы}

\begin{enumerate}
\item \textit{Динамика темпоральных графовых нейронных сетей непрерывного времени (CT-TGNN).} Впервые предложен метод, объединяющий передачу сообщений на уровне графа с динамикой нейронных обыкновенных дифференциальных уравнений на эволюционирующих топологиях. Многомасштабная темпоральная динамика с обучаемыми постоянными времени разрешает как быстрые переходные процессы, так и медленные структурные тенденции. Метод сопряжённых чувствительностей обеспечивает обучение с $O(1)$ памяти. Сертифицированные границы робастности выводятся из липшицевой гладкости управляющего векторного поля. Расширение SDE-TGNN достигает 99,0\% F1, ECE 0,042, сертифицированной состязательной робастности 76,2\% при радиусе возмущения 0,25 и кросс-доменного переноса 90,5\% F1 без дообучения.

\item \textit{Византийско-устойчивая федеративная оптимизация графов (FedLLM-API).} Разработан оригинальный федеративный алгоритм, совместно обеспечивающий устойчивую к повреждениям агрегацию (покоординатное усечённое среднее), формальную $(\varepsilon{=}0.85, \delta{=}10^{-5})$-дифференциальную приватность и выравнивание распределений по расстоянию Вассерштейна для графовых данных. Сходимость со скоростью $O(1/\sqrt{T})$ доказана даже при 30\% участников, отправляющих злонамеренно повреждённые обновления. Интеграция LLM обеспечивает 87,6\% F1 при обнаружении новых категорий атак в режиме zero-shot.

\item \textit{Многоуровневые темпоральные вложения графов (TripleE-TGNN).} Разработана архитектура тройного вложения, обучающая представления на уровнях сервисов, трассировок и узлов с помощью двойного темпорально-пространственного внимания и упругой консолидации весов. Архитектура достигает точности 96,8\% при снижении вычислительных затрат на 40\% за счёт структурированной разреженности.

\item \textit{Селективные модели пространства состояний для устойчивости к отравлению (MambaShield).} Разработан эффективный метод вывода на основе селективных слоёв пространства состояний со свёртками быстрого преобразования Фурье, снижающий сложность с $O(L^2)$ до $O(L \log L)$. Прогрессивная состязательная дистилляция робастности обеспечивает точность 91,4\% при 25\% отравления на этапе обучения. Предоставлены PAC-байесовские сертификаты обобщения.

\item \textit{Вывод с калиброванной неопределённостью (стохастический трансформер, UC-HGP).} Разработано вариационное байесовское внимание с оценкой неопределённости на иерархических гауссовских процессах, декомпозирующее предсказательную неопределённость на эпистемическую и алеаторную компоненты с ожидаемой ошибкой калибровки 0,078. Кросс-доменный перенос достигает обобщаемости 94,2\%.

\item \textit{Теоретико-игровая сертификация робастности.} Построена единая система сертификации, включающая равновесия Штакельберга, алгоритм мультипликативных весов с гарантиями отсутствия регрета ($R(T) \leq \sqrt{T \log |S|}$), распространение интервальных границ и рандомизированное сглаживание, адаптированное к дискретным графовым структурам. Регуляризация согласованности связывает все семь методов с доказанной сходимостью.

\item \textit{Предметно-ориентированная фундаментальная модель (CyberSecLLM).} Блоки селективного пространства состояний, дополненные ОДУ, достигают средней точности 89,1\% в режиме zero-shot на бенчмарке из 11 задач, что на 20,5 п.\,п.\ превышает GPT-4, с обработкой графовых последовательностей длиной в миллион токенов за 2,3 секунды.
\end{enumerate}

\subsection*{Теоретическая значимость}

Теоретическая значимость состоит в создании единой математической основы для робастного обучения на динамических графовых структурах. Эта основа объединяет пять областей математики, которые ранее не комбинировались для графовых данных:

\begin{enumerate}
\item Нейронные обыкновенные дифференциальные уравнения --- обеспечивающие сертифицированную динамику непрерывного времени с доказуемыми границами изменения выходов при возмущении входов.

\item Теория оптимального транспорта --- позволяющая выравнивать распределения данных между организациями при сохранении формальных гарантий приватности.

\item Селективные модели пространства состояний --- снижающие вычислительную сложность с квадратичной до близкой к линейной, делая обработку в реальном времени осуществимой на крупномасштабных графах.

\item Байесовский вывод --- декомпозирующий неопределённость предсказаний на сокращаемую (модельную) и несокращаемую (данных) компоненты, чтобы система знала, когда она не уверена.

\item Теоретико-игровая сертификация --- моделирующая взаимодействие защитника и противника как стратегическую игру с доказуемыми гарантиями робастности.
\end{enumerate}

Объединённая основа устанавливает шесть формальных гарантий --- скорости сходимости, сертификаты робастности по Липшицу, границы дифференциальной приватности, PAC-байесовские границы обобщения, сублинейный регрет и единая сходимость --- все подтверждены как теоретически (через доказательства), так и экспериментально (на 84,2 миллионах образцов).


\subsection*{Практическая значимость}

Практическая значимость определяется применимостью разработанных методов в пяти областях. В каждом случае указывается, \emph{кто} использует метод, \emph{как} он развёрнут и \emph{в какой области}:

\begin{enumerate}
\item \textbf{Здравоохранение.} Команды биомедицинского ИИ интегрируют CT-TGNN и SDE-TGNN в платформы поддержки клинических решений и лекарственного взаимодействия для обеспечения сертифицированно безопасных предсказаний на молекулярных графах и графах взаимодействия пациентов.

\item \textbf{Финансы.} Банки и финансовые регуляторы используют FedLLM-API для совместного обучения моделей обнаружения мошенничества на графах транзакций без обмена данными клиентов, обеспечивая приватность через гарантии дифференциальной приватности.

\item \textbf{Промышленный IoT.} Инженеры безопасности операционных технологий развёртывают MambaShield на периферийных шлюзах и контроллерах SCADA для обнаружения аномалий в реальном времени на графах датчиков, что обеспечивается его близкой к линейной вычислительной сложностью.

\item \textbf{Автономные системы.} Инженеры безопасности используют UC-HGP для получения калиброванных оценок уверенности для решений на графах слияния датчиков, позволяя системе маркировать неуверенные предсказания до того, как они приведут к небезопасным действиям.

\item \textbf{Кибербезопасность.} SOC-аналитики, инженеры кибербезопасности, исследователи MLSecOps и разработчики ИИ используют платформу с открытым исходным кодом RobustIDPS.ai (DOI:~10.5281/zenodo.19129512), которая интегрирует 13+ моделей нейронных сетей, четырёхуровневую систему защиты LLM с обнаружением 94,5\% по 517 сценариям атак и работает как плагин для Snort, Suricata и аналогичных СОПСВ.
\end{enumerate}

Практическая применимость подтверждена актами внедрения от трёх организаций: ООО <<Positive Technologies>> (федеративное обучение для коммерческих СОПСВ), Google Cloud Platform (периферийные вычисления, точность 96,8\%, снижение затрат на 40\%) и Suricata Open Source IDPS (адаптивное обнаружение с адаптацией к дрейфу концепта).

\subsection*{Достоверность научных результатов}

Достоверность подтверждается: строгой математической формализацией с доказательствами сходимости, устойчивости и сложности; экспериментальной валидацией на шести открытых наборах данных общим объёмом 84,2 миллиона размеченных образцов, покрывающих 84 уникальных типа атак; сравнительным анализом с современными базовыми линиями, демонстрирующим превосходство по ключевым метрикам; представлением на международных конференциях; публикациями в ведущих рецензируемых журналах (Scopus/Web of Science); а также актами внедрения от четырёх организаций.

\subsection*{Апробация результатов исследования}

Основные положения и результаты были представлены и обсуждены на следующих международных и российских научных конференциях:

\begin{enumerate}
\item Конференция молодых исследователей в области электротехники и электроники 2024 (ElCon-2024), Санкт-Петербург, Россия, IEEE;
\item Конференция молодых исследователей в области электротехники и электроники 2025 (ElCon-2025), ЭТУ-ЛЭТИ, Санкт-Петербург, Россия, IEEE;
\item Международная конференция <<Advances in Neural Computation, Machine Learning, and Cognitive Research IX>> (NEUROINFORMATICS-2025), Springer;
\item Третья Всероссийская научно-техническая конференция <<Кибернетика и информационная безопасность>> (КИБ-2025), НИЯУ МИФИ, Москва;
\item 2026 IEEE 27th International Conference of Young Professionals in Electron Devices and Materials (EDM-2026), Республика Алтай, Россия;
\item 11th IEEE Electronics System-Integration Technology Conference (IEEE ESTC 2026), Хельсинки, Финляндия;
\item IEEE Research and Technologies for Society and Industry (RTSI) 2026 Conference, Эспоо, Финляндия;
\item NeurIPS 2026: The 40th Annual Conference on Neural Information Processing Systems, Сидней, Австралия.
\end{enumerate}

\subsection*{Публикации}

Основные результаты опубликованы в 18 научных работах, в том числе: 5~статей в рецензируемых журналах, индексируемых в Scopus и Web of Science (Neurocomputing, Expert Systems with Applications, Complex and Intelligent Systems, Optical Memory and Neural Networks, IJASRE); 4~статьи в трудах международных конференций, индексируемых в Scopus и IEEE Xplore; 2~статьи приняты к публикации (Knowledge and Information Systems, Complex and Intelligent Systems); 7~статей на рецензировании в ведущих журналах IEEE Transactions и других изданиях.

\subsection*{Основные научные результаты, выносимые на защиту}

Результаты, выносимые на защиту, перечислены в порядке, отражающем их логическую взаимосвязь.

\begin{enumerate}
\item \textbf{Модели CT-TGNN и SDE-TGNN для сертифицированной динамики графов} --- оригинальные архитектуры непрерывного времени, в которых нейронные ОДУ и стохастические дифференциальные уравнения управляют эволюцией узлов и рёбер на динамических топологиях, обеспечивая сертифицированные границы робастности через анализ Липшица обученного векторного поля.

\item \textbf{Алгоритм FedLLM-API для распределённого обучения графов с сохранением приватности} --- федеративный оптимизатор с доказанной сходимостью, сочетающий византийско-устойчивую агрегацию усечённым средним с $(\varepsilon{=}0.85,\,\delta{=}10^{-5})$ дифференциальной приватностью и выравниванием оптимального транспорта Вассерштейна для гетерогенных фрагментов графов.

\item \textbf{Архитектура TripleE-TGNN для многоуровневого представления темпоральных графов} --- система тройного вложения, обучающая структуру на уровнях сервисов, трассировок и узлов с помощью двойного внимания и упругой консолидации весов, снижающая вычислительные затраты через структурированную разреженность и предотвращающая катастрофическое забывание.

\item \textbf{Архитектура MambaShield для эффективного вывода на графах} --- селективная модель пространства состояний, снижающая сложность вывода с $\mathcal{O}(L^{2})$ до $\mathcal{O}(L\log L)$ через свёртки БПФ, с прогрессивной состязательной дистилляцией и PAC-байесовскими сертификатами обобщения.

\item \textbf{Модели стохастического трансформера и UC-HGP для адаптации к дрейфу с учётом неопределённости} --- вариационное байесовское внимание с выводом на иерархических гауссовских процессах, обеспечивающее декомпозицию эпистемической и алеаторной неопределённости, онлайновое обнаружение дрейфа и гарантии сублинейного регрета при нестационарных графовых распределениях.

\item \textbf{Единая система сертификации, интегрирующая модели (1)--(5)} --- единая методология оценки, связывающая все разработанные модели через теоретико-игровые равновесия Штакельберга и регуляризацию согласованности, валидированная на 6 бенчмарк-наборах данных общим объёмом 84,2 миллиона размеченных образцов.
\end{enumerate}

\subsection*{Соответствие паспорту научной специальности}

Диссертация соответствует специальности 2.3.1 <<Системный анализ, управление и обработка информации, статистика>>:

\noindent П.4 <<Разработка методов и алгоритмов решения задач системного анализа, оптимизации, управления, принятия решений, обработки информации и искусственного интеллекта.>>

\noindent П.5 <<Разработка специального математического и алгоритмического обеспечения систем анализа, оптимизации, управления, принятия решений, обработки информации и искусственного интеллекта.>>

\subsection*{Структура и объём диссертации}

Диссертация состоит из введения, пяти глав, заключения, списка литературы из 190 наименований и трёх приложений. Основной текст составляет 210 страниц.

\section*{ОСНОВНОЕ СОДЕРЖАНИЕ РАБОТЫ}
\addcontentsline{toc}{section}{Основное содержание работы}

\textbf{Во Введении} обосновывается актуальность темы исследования, определяются объект, предмет, цели и задачи, научная новизна, теоретическая и практическая значимость, защищаемые положения, апробация результатов и структура диссертации.

\textbf{В Главе~1} представлен систематический анализ современных исследований в области робастного обучения на динамических графовых структурах. Выявлены четыре открытые проблемы, образующие логическую цепочку, определяющую структуру диссертации. Систематический анализ пробелов показывает, что ни один существующий подход не решает совместно задачи робастности, точности, эффективности и приватности при трёх рабочих условиях. Анализ охватывает состязательные возмущения и защиту графов, темпоральные и динамические графовые архитектуры, федеративное обучение на распределённых фрагментах графов, модели пространства состояний для эффективной темпоральной обработки, байесовское глубокое обучение и оценку неопределённости, теоретико-игровое моделирование взаимодействия защитника и противника, а также прямую совместимость с эволюционирующими протоколами. Выделены семь пробелов, каждый из которых сопоставлен конкретным пересечениям матрицы условие--требование (рис.~\ref{fig:main_arch}) и адресован методами М1--М7 (рис.~\ref{fig:venn_ar}). Анализ пробелов устанавливает мотивацию для решений, разработанных в Главе~2.

\textbf{В Главе~2} представлены разработанные методы и алгоритмы.

\textit{Метод М1: CT-TGNN --- Динамика темпоральных графовых нейронных сетей непрерывного времени.} Эволюция представлений узлов управляется параметризованными дифференциальными операторами на эволюционирующих топологиях графов:
\begin{equation}
\frac{dh_v(t)}{dt} = f_\theta(h_v(t), A(t), X(t), t)
\end{equation}
где $h_v(t)$ --- скрытое состояние узла $v$, $A(t)$ --- изменяющаяся во времени матрица смежности, $X(t)$ --- матрица признаков узлов.

Многомасштабная темпоральная динамика с постоянными времени $\tau_1, \ldots, \tau_K$, охватывающими несколько порядков величины, разрешает как быстрые переходные процессы, так и медленные структурные тенденции. Метод сопряжённых чувствительностей обеспечивает обучение с $O(1)$ памяти. Темпоральная адаптивная пакетная нормализация согласует непрерывную динамику со слоями нормализации.

\textbf{Теорема~1 (Сходимость):} Для константы Липшица $L_f$ и скорости обучения $\eta < 2/L_f$ обучение сходится со скоростью $O(1/\sqrt{T})$.

\textbf{Теорема~2 (Робастность):} Для возмущений входа $\|\delta X\| \leq \epsilon$ чувствительность выхода ограничена $\|\delta h(T)\| \leq \epsilon L_g T \exp(L_g T)$.

% ============================================================
% FIGURE 2: CT-TGNN Architecture
% ============================================================
На \textbf{рис.~\ref{fig:CT-TGNN}} представлена архитектура CT-TGNN для динамики темпоральных графов: метаданные зашифрованных потоков поступают через кодировщик признаков; три параллельные ветви нейронных ОДУ с различными постоянными времени ($\tau_{\text{fast}}, \tau_{\text{mid}}, \tau_{\text{slow}}$) захватывают динамику на уровне пакетов, сессий и кампаний; темпоральная адаптивная пакетная нормализация согласует непрерывную динамику; слияние масштабов с весами внимания формирует итоговое вложение узла для классификации.

\begin{figure}[H]
\centering
\includegraphics[width=0.65\textwidth]{figures/CT-TGNN_arch.pdf}
\vspace{-1.0cm}
\caption{Архитектура темпоральной графовой нейронной сети непрерывного времени CT-TGNN}
\label{fig:CT-TGNN}
\end{figure}

Стохастическое расширение SDE-TGNN на рис.~\ref{fig:SDE_TGNN} заменяет детерминированную динамику стохастическими дифференциальными уравнениями Ито и выводит аналитическое распространение моментов из уравнения Фоккера--Планка. Результаты: 99,0\% F1, ECE 0,042, 76,2\% сертифицированная состязательная робастность при $\varepsilon{=}0.25$, 90,5\% кросс-доменный перенос.
% ============================================================
% FIGURE: SDE-TGNN Architecture
% ============================================================
\begin{figure}[H]
\centering
\includegraphics[width=0.96\textwidth]{figures/SDE-TGNN_arch.pdf}
\vspace{-1cm}
\caption{Архитектура SDE-TGNN с распространением неопределённости по Фоккеру--Планку}
\label{fig:SDE_TGNN}
\end{figure}

На \textbf{рис.~\ref{fig:SDE_TGNN}} показана архитектура SDE-TGNN с распространением неопределённости по Фоккеру--Планку: темпоральный граф проходит через кодировщик для формирования начального состояния $h_0$. Два параллельных пути: решатель СДУ вычисляет стохастические траектории на нескольких масштабах ($\tau_s{=}0.5$\,с, $\tau_m{=}2$\,с, $\tau_l{=}10$\,с), а модуль Фоккера--Планка аналитически распространяет среднее $\mu_t$ и ковариацию $\Sigma_t$ через связанные ОДУ. Представление, объединённое вниманием, поступает в классификатор, формирующий предсказания и калиброванную неопределённость. Сертифицированный радиус получен из вероятностей сглаженного классификатора.

\textit{Метод М2: FedLLM-API --- Федеративная оптимизация графов с приватностью.} Агрегация использует покоординатное усечённое среднее с калиброванным гауссовским шумом:
\begin{equation}
\theta_{t+1} = \text{TrimMean}\left(\{\theta_i^t\}_{i=1}^N\right) + \mathcal{N}(0, \sigma^2 I)
\end{equation}

Выравнивание распределений через оптимальный транспорт:
\begin{equation}
W_2(\mu_i, \mu_j) = \inf_{\pi \in \Pi(\mu_i, \mu_j)} \left(\int \|x - y\|^2 \, d\pi(x,y)\right)^{1/2}
\end{equation}

\textbf{Теорема~3 (Сходимость при повреждениях):} Алгоритм сходится со скоростью $O(1/\sqrt{T})$ даже при 30\% повреждённых участников, сохраняя $(\varepsilon{=}0.85, \delta{=}10^{-5})$-дифференциальную приватность.

На \textbf{рис.~\ref{fig:FedLLM}} представлена архитектура FedLLM-API для совместного обучения на графах: синие пунктирные стрелки обозначают зашифрованные обновления градиентов с $(\varepsilon{=}0.85, \delta{=}10^{-5})$ дифференциальной приватностью; сервер применяет покоординатную агрегацию усечённым средним, устойчивую к 30\% злонамеренных участников; модуль ОТ Синкхорна выравнивает гетерогенные распределения клиентов; ветвь LLM обеспечивает классификацию zero-shot (87,6\% F1).
% ============================================================
% FIGURE 3: FedLLM-API Architecture
% ============================================================
\begin{figure}[H]
\centering \includegraphics[width=0.96\textwidth]{figures/FEDLLM_arch.pdf}
\vspace{-1cm}
\caption{Архитектура системы федеративного обучения FedLLM-API}
\label{fig:FedLLM}
\end{figure}

Интеграция LLM с рис.~\ref{fig:FedLLM} обеспечивает 87,6\% F1 при обнаружении новых аномальных категорий в режиме zero-shot.


\textit{Метод М3: TripleE-TGNN --- Многоуровневые темпоральные вложения графов.} Архитектура тройного вложения обучает представления на уровнях сервисов, трассировок и узлов с помощью двойного темпорально-пространственного внимания:
\begin{equation}
\alpha_{uv}^{(t)} = \frac{\exp(e_{uv}^{(t)})}{\sum_{w} \exp(e_{vw}^{(t)})}
\end{equation}

Упругая консолидация весов предотвращает катастрофическое забывание:
\begin{equation}
\mathcal{L}_{\text{EWC}} = \mathcal{L}_{\text{task}} + \frac{\lambda}{2} \sum_i F_i (\theta_i - \theta_i^*)^2
\end{equation}

Результаты: точность 96,8\%, снижение вычислительных затрат на 40\% за счёт структурированной разреженности.

% ============================================================
% FIGURE 4: TripleE-TGNN Architecture
% =========================================================

На \textbf{рис.~\ref{fig:TripleE_TGNN}} представлена архитектура TripleE-TGNN для многоуровневых представлений графов. Три параллельных кодировщика обрабатывают графы уровня сервисов, трассировок и узлов; двудольное межуровневое внимание интегрирует информацию по уровням абстракции; упругая консолидация весов предотвращает катастрофическое забывание; многозадачная голова одновременно классифицирует на всех трёх уровнях со штрафом согласованности по KL-дивергенции.

\begin{figure}[H]
\centering
\includegraphics[width=0.95\textwidth]{figures/TripleE_TGNN_arch.pdf}
\vspace{-1cm}
\caption{Архитектура многоуровневой темпоральной графовой сети TripleE-TGNN}
\label{fig:TripleE_TGNN}
\end{figure}

\textit{Метод М4: MambaShield --- Селективные модели пространства состояний.} Формулировка пространства состояний:
\begin{equation}
h_{t+1} = \bar{A} h_t + \bar{B} x_t, \quad y_t = C h_t + D x_t
\end{equation}

Свёртки БПФ снижают вычислительную сложность до $O(L \log L)$. Прогрессивная состязательная дистилляция робастности из ансамбля моделей-учителей:
\begin{equation}
\mathcal{L}_{\text{distill}} = \text{KL}\left(p_{\text{student}} \;\Big\|\; \frac{1}{K}\sum_{k} p_{\text{teacher}_k}\right)
\end{equation}

\textbf{Теорема~4 (PAC-байесовское обобщение):}
\begin{equation}
R(\theta) \leq \hat{R}(\theta) + \sqrt{\frac{\text{KL}(\rho \| \pi) + \ln(2N/\delta)}{2N}}
\end{equation}

Результаты: точность 91,4\% при 25\% отравления на этапе обучения.

Архитектура MambaShield для эволюционирующих ландшафтов угроз (\textbf{рис.~\ref{fig:mambashield}}): последовательные данные потоков проходят через блоки Mamba с учётом содержания, достигая линейной сложности; ансамбль из пяти моделей-учителей обеспечивает робастность через прогрессивную дистилляцию с учебным планом, постепенно увеличивающим долю отравления; PAC-байесовская регуляризация даёт формальный сертификат обобщения.
% ============================================================
% FIGURE 5: MambaShield Architecture
% ============================================================
\begin{figure}[H]
\centering
\includegraphics[width=0.65\textwidth]{figures/Mambashield_arch.pdf}
\vspace{-0.5cm}
\caption{Архитектура селективной модели пространства состояний MambaShield}
\label{fig:mambashield}
\end{figure}

\textit{Методы М5/М6: Стохастический трансформер и UC-HGP.} Вариационное байесовское внимание:
\begin{equation}
q_\phi(\alpha) = \mathcal{N}(\mu_\phi, \sigma_\phi^2)
\end{equation}

Нижняя граница свидетельства (ELBO):
\begin{equation}
\mathcal{L} = \mathbb{E}_{q_\phi}[\log p(y|x, \alpha)] - \text{KL}(q_\phi \| p)
\end{equation}

Иерархический гауссовский процесс: $f(x) \sim \mathcal{GP}(m(x), k(x, x'))$. Предсказательная неопределённость декомпозируется на эпистемическую и алеаторную компоненты через информационно-теоретическую основу: полная неопределённость равна предсказательной энтропии, разделяемой на алеаторную неопределённость (ожидаемая энтропия, несокращаемый шум данных) и эпистемическую неопределённость (взаимная информация между предсказаниями и параметрами модели, сокращаемая при наличии дополнительных данных). Система UC-HGP использует многомасштабное темпоральное моделирование через иерархию слоёв гауссовских процессов, работающих на временных масштабах уровня пакетов, сессий и кампаний, обеспечивая калиброванную оценку уверенности без накладных расходов на выборку Монте-Карло. Неправильно классифицированные образцы демонстрируют среднюю эпистемическую неопределённость 0,41 против 0,04 для правильных предсказаний, предоставляя сильный дискриминативный сигнал для автоматической сортировки оповещений. При 95\% покрытии риск селективного предсказания составляет лишь 0,88\%, а обнаружение OOD достигает 96,8\% AUROC для идентификации классов атак нулевого дня через пороговую обработку предсказательной энтропии. Результаты: ECE 0,078, кросс-доменный перенос 94,2\%, сокращение расследований на 43\%.

На \textbf{рис.~\ref{fig:UCHGP}} представлена архитектура UC-HGP: стохастическое внимание трансформера производит выборку матриц проекций из обученных гауссовских распределений; оценка Монте-Карло с пятьюдесятью прямыми проходами формирует предсказательные распределения; модуль декомпозиции неопределённости разделяет эпистемическую и алеаторную компоненты; температурное масштабирование обеспечивает калибровку; шлюз сортировки направляет предсказания с высокой неопределённостью аналитикам.

% ============================================================
% FIGURE: UC-HGP / Stochastic Transformer Architecture
% ============================================================
\begin{figure}[H]
\centering
\includegraphics[width=0.65\textwidth]{figures/UC_HGP_arch.pdf}
\vspace{-0.5cm}
\caption{Архитектура системы иерархических гауссовских процессов с калиброванной неопределённостью (UC-HGP)}
\label{fig:UCHGP}
\end{figure}

\textit{Метод М7: Теоретико-игровая сертификация робастности.} Равновесие Штакельберга с алгоритмом мультипликативных весов:
\begin{equation}
w_i^{(t+1)} = w_i^{(t)} \exp(-\eta \ell_i^{(t)})
\end{equation}

\textbf{Теорема~5 (Отсутствие регрета):} Кумулятивный регрет $R(T) \leq \sqrt{T \log |S|}$.

Система FedGTD моделирует безопасность мультиоблачных сред как стратегическую игру между множеством защитников и адаптивными противниками, учитывая крайнее разнообразие современных облачных сред. Система включает четыре ключевых нововведения: стратегическое теоретико-игровое моделирование с формальными гарантиями равновесия через теорему Какутани о неподвижной точке, византийско-устойчивое федеративное обучение, адаптированное к сильному дисбалансу данных (отношения дисбаланса $\rho_d \in [2.67, 99]$) и гетерогенным пространствам признаков, строгий мартингальный анализ сходимости в состязательных условиях с доказательством сходимости к $\varepsilon$-равновесию Нэша несмотря на невыпуклость целей и темпоральную нестационарность, а также предметно-адаптивное состязательное обучение, направленное на протокольные эксплойты (MQTT, Modbus), уязвимости контейнеров и корреляционные атаки. Оценка на 15,9 миллионах событий безопасности в трёх доменах демонстрирует точность обнаружения 95,7--96,9\% с приростом 4,8--8,8\% по сравнению с современными методами. Система сохраняет среднюю производительность 94,0\% при 20\% византийских повреждений (до 95,3\% против переворота меток), со снижением коммуникационных накладных расходов на 50,7\%.

Распространение интервальных границ, оценка Липшица на графовых операторах, рандомизированное сглаживание, адаптированное к дискретным графовым структурам. Регуляризация согласованности объединяет все семь методов с доказанной скоростью сходимости.

На \textbf{рис.~\ref{fig:FedGTD}} представлена теоретико-игровая система робастности для адаптивных графовых противников: гетерогенные данные из периферийных, контейнерных и распределённых сред поступают в предметно-ориентированные препроцессоры; модуль федеративного теоретико-игрового обучения вычисляет равновесия Штакельберга и Нэша через перечисление опор; алгоритм мультипликативных весов обеспечивает онлайн-адаптацию; эпистемическая неопределённость стохастического трансформера модулирует выигрыш защитника через полезность со штрафом за неопределённость.

% ============================================================
% FIGURE: FedGTD / Game-Theoretic Architecture
% ============================================================
\begin{figure}[H]
\centering
\includegraphics[width=0.65\textwidth]{figures/FedGTD_arch.pdf}
\vspace{-0.5cm}
\caption{Архитектура системы теоретико-игровой сертификации робастности (FedGTD)}
\label{fig:FedGTD}
\end{figure}

\textit{Фундаментальная модель CyberSecLLM.} Блоки селективного пространства состояний, дополненные ОДУ, достигают средней точности 89,1\% в режиме zero-shot на бенчмарке из 11 задач, что на 20,5 п.\,п.\ превышает GPT-4, с обработкой графовых последовательностей событий длиной в миллион токенов за 2,3 секунды. Как показано на \textbf{рис.~\ref{fig:CyberSecLLM}}, сетевые потоки графа токенизируются через четыре параллельных пути (внизу), обрабатываются через повторяющиеся группы из трёх блоков ODE-S6, блок кросс-внимания, привязанный к базе знаний разведки угроз $\mathcal{K}$, и слой прямого распространения MoE с предметно-специализированными экспертами. Задачные головы (вверху) включают голову точечного процесса для межсобытийного тайминга $\lambda^*(t)$.

\begin{figure}[H]
\centering \includegraphics[width=0.96\textwidth]{figures/CyberSecLLM_arch.pdf}
\vspace{-1cm}
\caption{Архитектурная основа CyberSecLLM}
\label{fig:CyberSecLLM}
\end{figure}


\textbf{В Главе~3} выполнена адаптация и согласование предметно-независимых методов к области кибербезопасности и системам обнаружения и предотвращения сетевых вторжений (СОПСВ). В главе определены специфические для кибербезопасности терминологии (поток, сессия, хост, сервис, PCAP, MITRE ATT\&CK, SOC), описано, как сырой сетевой трафик (файлы PCAP и данные CSV) преобразуется во временные графовые структуры через формальный конвейер данные--граф (узлы = IP-адреса с 12-мерными признаками хоста; рёбра = потоки с 83-мерными векторами признаков), представлена унифицированная таксономия классификации из 34 классов по 7 категориям атак, документирован конвейер инженерии 83 признаков, охватывающий статистику потоков, распределения межпакетных интервалов, флаговые индикаторы, статистику полезной нагрузки и производные отношения, а также указаны предметно-ориентированные гиперпараметры для каждого метода (CT-TGNN постоянные времени: 1/60/3600\,с; FedLLM-API бюджет ДП: $\varepsilon{=}0.85$; MambaShield ядро: 1024; стохастический трансформер: 50 выборок МК). Интегрированный конвейер обнаружения развёртывается как плагин для Snort и Suricata.

\textbf{В Главе~4} представлена экспериментальная валидация разработанных методов. Сетевой трафик моделируется как темпоральный граф $G_t = (V_t, E_t, X_t)$, где узлы представляют сетевые сущности, направленные рёбра кодируют коммуникационные потоки, а матрица признаков кодирует статистику трафика.

Экспериментальная оценка охватывает шесть бенчмарк-наборов данных, организованных в две взаимодополняющие группы, общим объёмом 84,2 миллиона размеченных образцов по 84 уникальным категориям атак:

\vspace{0.3cm}
\noindent\textit{Группа~1: Общий сетевой трафик} (IIS3D --- Интегрированная коллекция 3 наборов данных безопасности СОПСВ):

\begin{table}[H]
\centering
\caption{Бенчмарк-наборы данных, использованные в экспериментальной оценке}
\vspace{0.1cm}
\scriptsize
\begin{tabular}{@{}lccrcc@{}}
\toprule
\textbf{Набор данных} & \textbf{Домен} & \textbf{Год} & \textbf{Записи} & \textbf{Атаки} & \textbf{Призн.} \\
\midrule
\multicolumn{6}{@{}l}{\textit{Группа~1: Общий сетевой трафик (IIS3D)}} \\[2pt]
CIC-IoT-2023 & IoT & 2023 & 46,6M & 33 & 46 \\
CSE-CICIDS2018 & Корпор. & 2018 & 16,2M & 7 & 79 \\
UNSW-NB15 & Гибрид & 2015 & 2,5M & 9 & 49 \\
\midrule
\multicolumn{6}{@{}l}{\textit{Группа~2: Облако, микросервисы и периферия (ICS3D)}} \\[2pt]
Microsoft GUIDE & Облако & 2023 & 13,7M & 15 & 51 \\
Container Security & Микросерв. & 2023 & 3,2M & 8 & 93 \\
Edge-IIoT & Промышл. & 2022 & 2,0M & 12 & 69 \\
\midrule
\textbf{Всего} & & & \textbf{84,2M} & \textbf{84} & \\
\bottomrule
\end{tabular}
\end{table}

\vspace{-0.1cm}
Все модели используют унифицированный 46-мерный вход: наборы данных с большим числом признаков проходят отбор признаков по взаимной информации; наборы с меньшим числом признаков дополняются нулями. Предобработка включает нормализацию StandardScaler (RobustScaler для Container Security из-за разреженных признаков), построение $k$-NN графа ($k{=}10$) и стратифицированное темпоральное разделение (70\%/15\%/15\%). Обе группы наборов данных общедоступны на Kaggle.

Оценка использует единую систему из 23 метрик, охватывающую шесть категорий: качество обнаружения (точность, прецизионность, полнота, F1, MCC, макро-F1), качество калибровки (ECE, MCE, оценка Брайера, NLL), метрики неопределённости (AURC, OOD AUROC, покрытие), состязательная робастность (FGSM, PGD-20, сертифицированная точность при нескольких радиусах), вычислительная эффективность (пропускная способность, задержка, параметры) и приватность (расход бюджета дифференциальной приватности, деградация точности при $(\varepsilon, \delta)$-ДП).

% ============================================================
% FIGURE 7: Combined Experimental Results and Adversarial Robustness
% ============================================================
На \textbf{рис.~\ref{fig:roc_plot}} показана часть экспериментальной оценки и состязательной робастности:
(a)~ROC-кривые на CIC-IoT-2023: AUC-ROC 0,991 против 0,974 базовой линии. (b)~Диаграммы надёжности, сравнивающие ECE по методам. (c)~Робастная точность при PGD-атаках с увеличением бюджета возмущений. (d)~Сертифицированная точность как функция сертифицированного радиуса. (e)~Точность MambaShield при отравлении 0--40\%. (f)~Тепловая карта точности по наборам данных.

\begin{figure}[H]
\centering
\includegraphics[width=0.96\textwidth]{figures/ROC_6Datasets.pdf}

\caption{Сравнительная оценка и анализ состязательной робастности}
\label{fig:roc_plot}
\end{figure}


\textbf{Экспериментальные результаты:}

\begin{table}[H]
\centering
\caption{Производительность единой системы на бенчмарк-наборах данных}
\small
\begin{tabular}{lcccc}
\toprule
\textbf{Набор данных} & \textbf{Точн.} & \textbf{F1} & \textbf{Базов.} & \textbf{Улучш.} \\
\midrule
CIC-IoT-2023 & 94,7\% & 94,7\% & 91,8\% & +2,9 \\
CSE-CICIDS2018 & 93,9\% & 93,9\% & 89,7\% & +4,2 \\
UNSW-NB15 & 93,8\% & 93,8\% & 88,2\% & +5,6 \\
MS GUIDE & 92,1\% & 92,0\% & 87,5\% & +4,6 \\
Container Sec & 91,5\% & 91,4\% & 86,9\% & +4,6 \\
Edge-IIoT & 89,8\% & 89,7\% & 84,3\% & +5,5 \\
\bottomrule
\end{tabular}
\end{table}

\textbf{Абляционное исследование:} Удаление непрерывной динамики: $-6{,}8\%$; многоуровневых вложений: $-8{,}3\%$; робастной агрегации: $-21{,}4\%$; оптимального транспорта: $-15{,}8\%$; компонента LLM: $-53{,}4$ F1; дистилляции: $-9{,}7\%$; вариационного внимания: $-114$ ECE.

\textbf{Состязательная робастность:} CT-TGNN сохраняет точность 76,2\% при $\varepsilon{=}0.25$; единая система достигает 89,6\% при комбинированных возмущениях.

\textbf{Кросс-доменный перенос:} LibriSpeech 94,2\% F1, MIMIC-III/eICU 89,7\%.

\textbf{Вычислительная производительность:} 8,7 миллиона событий/сек, 47\,мс задержка.

\textbf{Приватность:} $(\varepsilon{=}0.85, \delta{=}10^{-5})$, деградация точности лишь 3,1\%.

\textbf{Оценка неопределённости:} ECE 0,078, сокращение расследований на 43\%.


\textbf{В Главе~5} описана программная платформа RobustIDPS.ai (DOI:~10.5281/zenodo.19129512), реализующая все разработанные методы как единую развёртываемую систему (рис.~\ref{fig:Main_arch_v1b}). Платформа содержит 46 страниц в 9 функциональных группах (Командный центр ИИ, Активная защита ИИ, SOC-аналитика, Поверхности атак LLM, Стандарты MLSecOps, Безопасность и управление ИИ, Новые методы ИИ, Данные и модели ИИ, Система), интегрируя 13+ моделей нейронных сетей. Пакет SOC-аналитики обеспечивает агентное авторасследование, поиск угроз на естественном языке и автоматическую генерацию правил СОВ для Snort/Suricata. Четырёхуровневая система защиты LLM достигает обнаружения 94,5\% по 517 сценариям атак с четырьмя коммерческими бэкендами (Claude, GPT-4o, Gemini, DeepSeek). Модуль Multi-Agent PQC-IDS расширяет обнаружение на постквантовые протоколы. Пользовательское исследование подтвердило 34\% сокращение времени сортировки аналитиками. Акты внедрения от четырёх организаций: Росатом (95,7\%), Positive Technologies, GCP (96,8\%, снижение затрат на 40\%) и Suricata IDPS.
\begin{figure}
\centering
\includegraphics[width=\textwidth]{figures/Main_arch_v1b.pdf}
\caption{Архитектурная основа RobustIDPS}
\label{fig:Main_arch_v1b}
\end{figure}

\section*{ОСНОВНЫЕ РЕЗУЛЬТАТЫ И ВЫВОДЫ}
\addcontentsline{toc}{section}{Основные результаты и выводы}

\begin{enumerate}
\item Разработана математическая модель темпоральных графовых нейронных сетей непрерывного времени (CT-TGNN) с многомасштабной динамикой и сертифицированной робастностью по Липшицу. SDE-TGNN: 99,0\% F1, ECE 0,042, 76,2\% сертифицированная робастность при $\varepsilon{=}0.25$, 90,5\% кросс-доменный перенос.

\item Создан федеративный алгоритм (FedLLM-API) с робастной агрегацией, $(\varepsilon, \delta)$-дифференциальной приватностью и выравниванием оптимального транспорта. Сходимость $O(1/\sqrt{T})$ при 40\% повреждённых участников. LLM: 87,6\% F1 при обнаружении новых категорий.

\item Разработаны многоуровневые вложения (TripleE-TGNN) с двойным вниманием и упругой консолидацией весов. Точность 96,8\%, снижение затрат на 40\%.

\item Создан метод селективного пространства состояний (MambaShield) со сложностью $O(L \log L)$, достигающий точности 91,4\% при 25\% отравления, с PAC-байесовскими сертификатами.

\item Разработана интеграция вариационного внимания с гауссовскими процессами для декомпозиции неопределённости. ECE 0,078, сокращение расследований на 43\%.

\item Построена теоретико-игровая сертификация с равновесиями Штакельберга, онлайн-обучением без регрета и объединением методов через регуляризацию согласованности.

\item Создана фундаментальная модель CyberSecLLM: точность 89,1\% (улучшение на 20,5 п.\,п.\ над GPT-4), миллион токенов за 2,3\,с.

\item Валидация на 84,2 миллионах образцов, 84 категориях. Единая система: 94,7\% (CIC-IoT), 93,9\% (CICIDS), 93,8\% (UNSW), превышая базовые линии на 2,9--5,6 п.\,п.

\item Реализована платформа RobustIDPS.ai (DOI:~10.5281/zenodo.19129512), 8 моделей, четырёхуровневая защита LLM, обнаружение 94,5\% по 517 сценариям.

\item Внедрение в четырёх организациях: Росатом (95,7\%), Positive Technologies (федеративное обучение), GCP (96,8\%, снижение затрат на 40\%), Suricata (адаптация к дрейфу). Акты получены от всех организаций.
\end{enumerate}

Система обеспечивает одновременное решение задач состязательной робастности, приватности, обобщаемости и эффективности на единой математической основе, объединяющей динамику графов непрерывного времени, оптимальный транспорт, байесовский вывод, модели пространства состояний и теорию игр.

\vspace{1cm}

\section*{ПУБЛИКАЦИИ ПО ТЕМЕ ДИССЕРТАЦИИ}
\addcontentsline{toc}{section}{Публикации по теме диссертации}

\sloppy

\subsection*{Статьи в журналах, индексируемых в Scopus/Web of Science}

\begin{enumerate}
\item \textit{Anaedevha R.N., Trofimov A.G., Borodachev Y.V.} Uncertainty-Calibrated Hierarchical Gaussian Processes for Intrusion Detection with Multi-Scale Temporal Modeling // Neurocomputing. 2026. V.677. P.133105. DOI:~10.1016/j.neucom.2026.133105 \textbf{(Q1)}

\item \textit{Anaedevha R.N., Trofimov A.G., Borodachev Y.V.} MambaShield: Selective State-Space Models for Poisoning-Resilient Sequential Modeling // Expert Systems with Applications. 2026. P.131175. DOI:~10.1016/j.eswa.2026.131175 \textbf{(Q1)}

\item \textit{Anaedevha R.N., Trofimov A.G., Borodachev Y.V.} Temporal Adaptive Neural Ordinary Differential Equations with Deep Spatio-Temporal Point Processes for Real-Time Network Intrusion Detection // Complex and Intelligent Systems. Springer. 2026. DOI:~10.1007/s40747-026-02291-7 \textbf{(Q1)}

\item \textit{Anaedevha R.N., Trofimov A.G.} Improved Robust Adversarial Model against Evasion Attacks on Intrusion Detection Systems // Optical Memory and Neural Networks. 2024. V.33(Suppl.3). Pp.S414--S423. DOI:~10.3103/S1060992X24700681 \textbf{(Q3)}

\item \textit{Anaedevha R.N. et al.} Cyber Security Framework for Nigerian Civil Aviation Authority // International Journal of Advances in Scientific Research and Engineering. 2020. V.6(1). Pp.188--195. DOI:~10.31695/IJASRE.2020.33695
\end{enumerate}

\subsection*{Труды конференций}

\begin{enumerate}
\setcounter{enumi}{5}
\item \textit{Anaedevha R.N., Trofimov A.G.} Integrated Deep Q-Networks and Actor-Critic Models for Enhanced Poisoning Attack Detection // 2025 Conf.\ Young Researchers in Electrical and Electronic Engineering (ElCon). IEEE Xplore. 2025.

\item \textit{Anaedevha R.N.} Application of Machine Learning Models for Device Identification in Wireless Network Traffic // 2024 ElCon. IEEE Xplore. 2024. Pp.104--110. DOI:~10.1109/ElCon61730.2024.10468413

\item \textit{Anaedevha R.N., Trofimov A.G.} Stochastic Multimodal Transformer with Uncertainty Quantification for Robust Network Intrusion Detection // Advances in Neural Computation, Machine Learning, and Cognitive Research IX. NEUROINFORMATICS 2025. Springer. DOI:~10.1007/978-3-032-07690-8\_35

\item \textit{Anaedevha R.N. et al.} Sustainable Business Model for Airline Operation in a Developing Economy // Proc.\ Intl.\ Conf.\ Science, Technology, Education, Arts, Management. 2018. Pp.167--174.

\item \textit{Anaedevha R.N., Trofimov A.G.} SDE-TGNN: Multi-Scale Temporal Graph Neural Network with Stochastic Dynamics for Cloud and Industrial Network Intrusion Detection // Proc.\ 27th IEEE Int.\ Conf.\ Young Professionals in Electron Devices and Materials (EDM-2026), Altai Republic, Russia. IEEE Xplore. 2026.
\end{enumerate}

\subsection*{Принятые к публикации}

\begin{enumerate}
\setcounter{enumi}{9}
\item \textit{Anaedevha R.N., Trofimov A.G., Borodachev Y.V.} Stochastic PAC-Bayesian Transformers for Network Intrusion Detection and NLP Applications // Knowledge and Information Systems. Springer. 2025. Принята.

\item \textit{Anaedevha R.N., Trofimov A.G., Borodachev Y.V.} Byzantine-Resilient Stochastic Games for Federated Multi-Cloud Intrusion Detection // Complex and Intelligent Systems. Springer. 2025. Принята.
\end{enumerate}

\subsection*{На рецензировании}

\begin{enumerate}
\setcounter{enumi}{11}
\item \textit{Anaedevha R.N., Trofimov A.G.} TripleE-TGNN: Triple-Embedding Temporal Graph Neural Networks for Multi-Granularity Microservices Security // IEEE Trans.\ Dependable and Secure Computing. 2026.

\item \textit{Anaedevha R.N., Trofimov A.G., Borodachev Y.V.} PQ-IDPS: Adversarially Robust Intrusion Detection for Post-Quantum Encrypted Traffic // IEEE Trans.\ Neural Networks and Learning Systems. 2026.

\item \textit{Anaedevha R.N., Trofimov A.G., Borodachev Y.V.} FedLLM-API: Privacy-Preserving LLM Federation for Zero-Day API Threat Detection // IEEE Trans.\ Information Forensics and Security. 2026.

\item \textit{Anaedevha R.N., Trofimov A.G., Borodachev Y.V.} CT-TGNN: Continuous-Time Temporal Graph Neural Networks for Real-Time Threat Propagation // IEEE Trans.\ Services Computing. 2026.

\item \textit{Anaedevha R.N., Trofimov A.G., Borodachev Y.V.} Hybrid Spatial-Temporal Deep Learning for Privacy-Preserving Encrypted Traffic Intrusion Detection // IEEE Access. 2026.

\item \textit{Anaedevha R.N., Trofimov A.G., Borodachev Y.V.} Differentially Private Optimal Transport for Multi-Cloud Intrusion Detection // IEEE Access. 2026.

\item \textit{Taremwa M., Anaedevha R.N., Trofimov A.G.} Robust Audio-Image Steganography using Cross-Modal Transformer Models // Journal of Computer Virology and Hacking Techniques. Springer. 2024.
\end{enumerate}

\end{document}
